\section{Robust information-aware active observation design}
\label{sec:active-information-v15}

We extend a passive bistatic link $(i,j,q)$ to an active observation
$a=(i,j,q,m)$, where the acquisition mode $m$ determines the sensing-power
scale $p_m$, the number of independent looks $L_m$, and whether refined
delay--Doppler processing is used.  The central design principle is to price
an observation by the hypothesis-separation information that reaches the
fusion node, while reserving final detection claims for exact-likelihood
Monte Carlo evaluation.

For $L$ independent complex samples, we use
\begin{equation}
 H_0:y_\ell\sim\mathcal{CN}(0,\sigma^2),\qquad
 H_1:y_\ell\sim\mathcal{CN}(0,(1+\gamma)\sigma^2).
\end{equation}
The directed divergences and their Jeffreys sum are
\begin{align}
 D(P_1\Vert P_0)&=L\left[\gamma-\log(1+\gamma)\right],\\
 D(P_0\Vert P_1)&=L\left[\log(1+\gamma)-\frac{\gamma}{1+\gamma}\right],\\
 J(P_1,P_0)&=L\frac{\gamma^2}{1+\gamma}.
\end{align}
The final expression is exactly the separation of the centered LLR means.
If report arrival is observed, independent of the hypothesis, and succeeds
with probability $\chi_a$, the received divergence is exactly $\chi_a$ times
the local divergence.  Conditional independence then makes the information
of a bundle additive.

We represent target-aspect uncertainty by a finite, path-dependent scenario
set.  For viewing-bisector azimuth $\phi_{ijq}$ and scenario angle $\theta_s$,
the transparent abstraction
\begin{equation}
 g_s(\phi)=g_{\min}+(1-g_{\min})\cos^2(\theta_s-\phi)
\end{equation}
gives $\gamma_{a,s}=p_m g_s(\phi_{ijq})\gamma_{ijq}^{\rm coarse/fine}$.
Unlike a scalar RCS interval, these scenarios can rank different bistatic
views differently.

For one target--fusion pair, the information-pricing problem is
\begin{equation}
 \max_x\;\min_{s\in\mathcal S}\sum_a x_a I_{a,s}
 -\lambda_E\sum_a x_a e_a-\lambda_R\sum_a x_a r_a,
\end{equation}
subject to one mode per physical link, observation-count, energy, and
report-feasibility constraints.  A depth-first branch-and-bound oracle uses a
scenario-wise top-remaining-information relaxation and ignores future
nonnegative costs.  This is a valid optimistic bound.  Exhaustion therefore
certifies optimality over the declared shortlist; node-limit truncation is
reported through the incumbent--upper-bound gap.

Each selected observation is subsequently evaluated by a mixed-mode exact-LLR
detector, in which both $\gamma_a$ and $L_a$ are observation specific.  A
strict threshold comparison preserves the intended false-alarm semantics in
the presence of the erasure atom at zero.  Refinement incurs
$C_a=L_a(C_{\rm MF}+C_{\rm LLR})+\mathbf{1}_{\rm refined}C_{\rm DD}$ cycles.

To expose the power externality without making the master nonlinear, active
columns are evaluated under a full-load interference envelope: every
interfering UAV radiates at the largest declared sensing-power scale, while
the desired observation uses its selected scale.  A column stores
scenario-wise $P_D$, received information, energy, CPU, receiver load and
report use.  The fleet-level MILP lexicographically minimizes the worst and
total detection deficits, followed by energy, reports and computation, under
shared UAV and fusion-node budgets.  Its certificate applies to the generated
column pool.

Across 90 paired target instances from nine geometry clusters, active modes
increased worst-scenario $P_D$ by $0.03025$ on average (95\% cluster-bootstrap
CI, $[0.01998,0.04006]$) while mean scenario $P_{FA}$ remained near $0.05$.
In a separate nine-instance, three-UAV/two-target global diagnostic, an
independent 32,768-sample holdout replay showed that active columns increased
mean-target $P_D$ by $0.01717$ (95\% CI, $[0.00360,0.03571]$) under shared
energy, CPU, reporting and load constraints.  The worst-target gain was
$0.01282$, but its interval $[-0.00066,0.03345]$ crossed zero.  These results
support an average detection consequence under the declared conditions; they
do not establish a worst-target improvement or a full-scale promotion result.
